\section{Die Multi-Level-Monte-Carlo-Methode}
Die Multi-Level-Monte-Carlo-Methode (MLMC) beruht genauso wie die klassische Monte-Carlo-Methode auf wiederholtem zuf"alligem Sampling. Jedoch werden diese Samples nun auf unterschiedlichen Level von unterschiedlicher Genauigkeit ausgef"uhrt.
Die MLMC-Methode kann die Rechenkosten der klassischen MC-Methode stark reduzieren, indem man die meisten Samples auf einem Level mit niedriger Genauigkeit und entsprechend niedrigen Kosten nimmt und nur wenige Samples bei hoher Genauigkeit und entsprechend hohen Kosten.\\
Auch bei der MLMC-Methode ist das Ziel, die Erwartung $\mathbb{E}(X)$ einer Zufallsvariable $X$ zu approximieren. Dazu nutzt man  eine Folge von Approximationen $X_0,X_1,...,X_L$ mit ansteigender Genauigkeit und Kosten, die f"ur $L\longrightarrow\infty$ gegen $X$ konvergieren.\\
Die Grundlage der Multi-Level-Methode ist die Teleskopsumme. Genauer:
\begin{align*}
\mathbb{E}(X_L)=\mathbb{E}(X_0)+\sum_{l=1}^L \mathbb{E}(X_l-X_{l-1}).
\end{align*}
Durch die Linearit"at des Erwartungswertes ist diese Gleichung trivialerweise erf"ullt. Dabei wird jede Erwartung $\mathbb{E}(X_l-X_{l-1})$ durch die Monte-Carlo-Methode approximiert. Man muss jedoch beachten, dass ein Sample von $X_l-X_{l-1}$ auf dem Level $l$ sowohl eine Simulation von $X_{l-1}$ als auch von $X_l$ ben"otigt.\\
Die MLMC-Methode funktioniert, falls die Varianz $\text{Var}[X_l-X_{l-1}]$ f"ur $l\rightarrow\infty$ gegen $0$ konvergiert. Dies ist jedoch der Fall, falls $X_l$ und $X_{l-1}$ dieselbe Zufallsvariable approximieren.\vspace{1cm}

F"ur den MLMC-Algorithmus angewandt auf unser Nadelproblem definieren wir $D_0=Y_0$ und f"ur $l=1,2,...$
\begin{align*}
D_l(x,\alpha):=(Y_l-Y_{l-1})(x,\alpha)=\sum_{j=1}^{m_l} f(x,\alpha;\lambda_j^l)\mathbb{I}_{(\lambda_{j-1}^l,\lambda_j^l]}-\sum_{j=1}^{m_{l-1}}f(x,\alpha;\lambda_j^{l-1})\mathbb{I}_{(\lambda_{j-1}^{l-1},\lambda_j^{l-1}]}
\end{align*}
Zuerst betrachten wir die Varianz und den Erwartungswert von $D_l$.
\newpage
\begin{prop}\label{k}
Es gilt $\mathbb{E}(D_0)=a(1)\mathbb{I}_{(0,1]}$ und f"ur jedes $l=1,2,...$ gilt
\begin{align*}
\mathbb{E}(D_l)=\sum_{j=1}^{m_{l-1}}(a(\lambda_{2j-1}^l)-a(\lambda_{2j}^l))\mathbb{I}_{(\lambda_{2j-2}^l,\lambda_{2j-1}^l]}.
\end{align*}
Zus"atzlich gilt f"ur alle $l\in\mathbb{N}$
\begin{align*}
\|D_l-\mathbb{E}[D_l]\|^2_{L^2(\Omega,H)}=\frac{1}{\pi m_l}\left(1-\frac{2}{\pi m_l}\right).
\end{align*}
\end{prop}
\begin{proof}
Es gilt $\lambda_{j-1}^{l-1}=\frac{j-1}{2^{l-1}}=\frac{2(j-1)}{2^l}=\lambda^l_{2(j-1)}$ und analog $\lambda_j^{l-1}=\lambda_{2j}^l$.\\
Deshalb gilt auch 
\begin{align*}
\mathbb{I}_{(\lambda_{j-1}^{l-1},\lambda_j^{l-1}]}=\mathbb{I}_{(\lambda_{2(j-1)}^{l},\lambda_{2j}^{l}]}=\mathbb{I}_{(\lambda_{2(j-1)}^{l},\lambda_{2j-1}^{l}]}+\mathbb{I}_{(\lambda_{2j-1}^{l},\lambda_{2j}^{l}]} 
\end{align*}
f"ur alle $ l=1,2,...$ und $ j=1,...,m_l$.\\
Damit folgt
\begin{align*}
\mathbb{E}[D_l]=&\mathbb{E}\left(\sum_{j=1}^{m_l}f(x,\alpha;\lambda_j^l)\mathbb{I}_{(\lambda_{j-1}^{l},\lambda_j^{l}]}\right)-\mathbb{E}\left(\sum_{j=1}^{m_{l-1}}f(x,\alpha;\lambda_j^{l-1})\mathbb{I}_{(\lambda_{j-1}^{l-1},\lambda_j^{l-1}]}\right)\\=&\sum_{j=1}^{m_l}a(\lambda_j^l)\mathbb{I}_{(\lambda_{j-1}^{l},\lambda_j^{l}]}-\sum_{j=1}^{m_{l-1}}a(\lambda_j^{l-1})\mathbb{I}_{(\lambda_{j-1}^{l-1},\lambda_j^{l-1}]}\\=&\sum_{j=1}^{m_{l-1}}\left(a(\lambda_{2j-1}^l)\mathbb{I}_{(\lambda_{2j-2}^{l},\lambda_{2j-1}^{l}]}+a(\lambda_{2j}^l)\mathbb{I}_{(\lambda_{2j-1}^{l},\lambda_{2j}^{l}]}\right)\\-&\sum_{j=1}^{m_{l-1}}\left(a(\lambda_{j}^{l-1})\mathbb{I}_{(\lambda_{2(j-1)}^{l},\lambda_{2j-1}^{l}]}+a(\lambda_{j}^{l-1})\mathbb{I}_{(\lambda_{2j-1}^{l},\lambda_{2j}^{l}]}\right)\\=& \sum_{j=1}^{m_{l-1}}a(\lambda_{2j-1}^l)-a(\lambda_{2j}^l))\mathbb{I}_{(\lambda^l_{2j-2},\lambda^l_{2j-1}]}=\sum_{j=1}^{m_l}\left(-\frac{2}{\pi m_l}\right)\mathbb{I}_{(\lambda^l_{2j-2},\lambda^l_{2j-1}]}.
\end{align*}
Dies beweist den ersten Teil der Proposition. F"ur den zweiten Teil gilt
\allowdisplaybreaks
\begin{align*}
\|D_l-\mathbb{E}(D_l)\|^2_{L^2(\Omega,H)}
=&\mathbb{E}
\int_0^1|\sum_{j=1}^{m_l}
f(\cdot;\lambda_j^l)\mathbb{I}_{(\lambda_{j-1}^l,\lambda_j^l]}
\sum_{j=1}^{m_{l-1}}f(\cdot;\lambda_j^{l-1})\mathbb{I}_{(\lambda_{j-1}^{l-1},\lambda_j^{l-1}]}\\
-&\sum_{j=1}^{m_{l-1}}\left(-\frac{2}{\pi m_l}\right)\mathbb{I}_{(\lambda^l_{2j-2},\lambda^l_{2j-1}]}|^2\td{\lambda}\\
=&\mathbb{E}\int_0^1|\sum_{j=1}^{m_{l-1}}f(\cdot;\lambda_{2j-1}^{l})-f(\cdot;\lambda_{2j}^{l})+\frac{2}{\pi m_l})\mathbb{I}_{(\lambda^l_{2j-2},\lambda^l_{2j-1}]}|^2\td{\lambda}\\
=& \sum_{j=1}^{m_{l-1}}\frac{1}{m_l}\left(\mathbb{E}|f(\cdot;\lambda^l_{2j-1})-f(\cdot;\lambda_{2j}^l)|^2-\left(\frac{2}{\pi m_l}\right)^2\right)\\
=&\frac{1}{m_l}\sum_{j=1}^{m_{l-1}}\frac{2}{\pi m_l}\left(1-\frac{2}{\pi m_l}\right)=\frac{1}{m_l}\left(m_{l-1}\left(\frac{2}{\pi m_l}\left(1-\frac{2}{\pi m_l}\right)\right)\right)
\\=&
\frac{1}{\pi m_l}\left(1-\frac{2}{\pi m_l}\right),
\end{align*}
da $\frac{m_{l-1}}{m_l}=\frac{1}{2}$.
\end{proof}

